\documentclass[../../main.tex]{subfiles}
\begin{document}

{\bf [解]}
題意の曲線を$C: y=f(x)=x^4-6x^2$とおく．$f'(x)=4x^3-12x$より，$x=t$における$C$の接線$\ell$は
\begin{align}
\ell: y=(4t^3-12t)x-3t^4+6t^2 \label{eq:1}
\end{align}
と書ける．
これが$P(\alpha,\beta)$を通る時，
\begin{align}
\beta=(4t^3-12t)\alpha-3t^4+6t^2 \quad \cdots \text{①} \label{eq:2}
\end{align}
を満たす．
$C$の2重接線は$y=-9$だけであり，題意から$P$が$y>x^4-6x^2$を動く限りは考えなくて良い．
従って題意の条件は\cref{eq:2}が$t$に関して4実解を持てば良い．
そこで，\cref{eq:2}を整理して
\begin{align}
 g(t)=3t^4-4\alpha t^3-6t^2+12\alpha t+\beta \label{eq:3}
\end{align}
とおくと\cref{eq:2}は$g(t)=0$と同じである．$g(t)$の挙動を調べるため一階微分を計算すると
\begin{align}
 g'(t)=12t^3-12\alpha t^2-12t+12\alpha=12(t-\alpha)(t+1)(t-1)
\end{align}
より，$\alpha$の値によって以下のような増減表を得る．


表を得る．まず，$t=0$が解でないので$\beta\ne0$．

\subparagraph{1. $\alpha\le-1$}

この時の増減表は以下の通り．

\begin{table}[htb]
\begin{center}
\caption{$g(t)$の増減表}
\begin{tabular}{c|ccccccc}
$t$ & & $\alpha$ & & $-1$ & & $1$ & \\
\hline
$g'$ & $-$ & $0$ & $+$ & $0$ & $-$ & $0$ & $+$ \\
\hline
$g$ & $\searrow$ & & $\nearrow$ & & $\searrow$ & & $\nearrow$
\end{tabular}
\end{center}
\end{table}

従って題意を満たす条件は
\begin{align}
g(\alpha)<0,\ g(1)<0,\ g(-1)>0 \label{eq:4}
\end{align}
である．

\subparagraph{2. $-1\le\alpha\le 1$}

この時の増減表は以下の通り．

\begin{table}[htb]
\begin{center}
\caption{$g(t)$の増減表}
\begin{tabular}{c|ccccccc}
$t$ & & $-1$ & & $\alpha$ & & $1$ & \\
\hline
$g'$ & $-$ & $0$ & $+$ & $0$ & $-$ & $0$ & $+$ \\
\hline
$g$ & $\searrow$ & & $\nearrow$ & & $\searrow$ & & $\nearrow$
\end{tabular}
\end{center}
\end{table}

従って題意を満たす条件は
\begin{align}
 g(-1)<0,\ g(1)<0,\ g(\alpha)>0 \label{eq:5}
\end{align}
である．

 \subparagraph{3. $1\le\alpha$}

この時の増減表は以下の通り．

\begin{table}[htb]
\begin{center}
\caption{$g(t)$の増減表}
\begin{tabular}{c|ccccccc}
$t$ & & $-1$ & & $1$ & & $\alpha$ & \\
\hline
$g'$ & $-$ & $0$ & $+$ & $0$ & $-$ & $0$ & $+$ \\
\hline
$g$ & $\searrow$ & & $\nearrow$ & & $\searrow$ & & $\nearrow$
\end{tabular}
\end{center}
 \end{table}

従って題意を満たす条件は
\begin{align}
 g(-1)<0,\ g(\alpha)<0,\ g(1)>0 \label{eq:6}
\end{align}
である．
\casesend

以上の場合分けより，$g(1), g(-1), g(\alpha)$の符号を調べれば良い．\cref{eq:3}より
\begin{align}
g(1)=\beta+8\alpha-3,\qquad g(-1)=\beta-8\alpha-3,\qquad g(\alpha)=-\alpha^4+6\alpha^2+\beta
\end{align}
である．
ここで$g(\alpha)=\beta-f(\alpha)$であり，題意より$\beta>f(\alpha)$を満たすから常に$g(\alpha)>0$である．
よって\cref{eq:4,eq:5,eq:6}より$g(\alpha)<0$を要求する $1^\circ,3^\circ$ は起こりえず $2^\circ$（$-1\le\alpha\le1$）のみが可能である．
よって求める条件は\cref{eq:5}で，$g(\alpha)>0$が自動的にみたされることから
\begin{align}
g(-1)<0\ \text{かつ}\ g(1)<0 \iff \beta<3-8\alpha\ \text{かつ}\ \beta<3+8\alpha
\end{align}
に帰着する．以上から，もとめる領域$D$は
\begin{align}
D=\left\{(\alpha,\beta)\ \middle|\ -1<\alpha<1,\ \ \alpha^4-6\alpha^2<\beta<3-8|\alpha|\right\}
\end{align}
であり，図示すると下図（境界含まず）である．

\begin{figure}[htb]
\begin{center}
\begin{tikzpicture}[xscale=1.2, yscale=0.5, >=stealth, every node/.style={font=\small}]

  % 1. 領域 D の斜線塗りつぶし
  % (0,3) -> (1,-5) -> [曲線 y=x^4-6x^2] -> (-1,-5) -> (0,3)
  \fill[pattern=north east lines, pattern color=black!50] 
    (0,3) -- (1,-5) 
    -- plot[domain=1:-1, smooth] (\x, {\x*\x*\x*\x - 6*\x*\x}) 
    -- cycle;

  % 2. 補助線（点線）
  % x = 1, -1 から x軸・y軸(-5) への垂線
  \draw[dashed, gray!70] (1, 0) -- (1, -5) -- (-1, -5) -- (-1, 0);
  \draw[dashed, gray!70] (0, -5) -- (1, -5);
  % y = -9 (極小値) の補助線
  \draw[dashed, gray!70] (-1.732, -9) -- (1.732, -9);
  \draw[dashed, gray!70] (-1.732, 0) -- (-1.732, -9);
  \draw[dashed, gray!70] (1.732, 0) -- (1.732, -9);

  % 3. 座標軸
  \draw[->, thick] (-3.2, 0) -- (3.2, 0) node[right] {$\alpha$};
  \draw[->, thick] (0, -9.8) -- (0, 4.0) node[above] {$\beta$};
  \node[below left] at (0,0) {$O$};

  % 4. 曲線と直線 (境界線)
  % 全体としての曲線 C: y = x^4 - 6x^2
  \draw[domain=-2.55:2.55, smooth, thick, blue!80!black] plot (\x, {\x*\x*\x*\x - 6*\x*\x}) node[right] {$C$};
  
  % 直線 y = 3 - 8|x| (領域の外側への延長部)
  \draw[thick, gray!60] (1, -5) -- (1.2, -6.6);
  \draw[thick, gray!60] (-1, -5) -- (-1.2, -6.6);

  % 領域 D の境界線（「境界を含まない」ため領域該当部を破線で強調）
  \draw[thick, dashed, red!80!black] (0,3) -- (1,-5);
  \draw[thick, dashed, red!80!black] (0,3) -- (-1,-5);
  \draw[thick, dashed, red!80!black, domain=-1:1, smooth] plot (\x, {\x*\x*\x*\x - 6*\x*\x});

  % 5. 軸上の目盛りテキスト
  \node[left] at (0, 3) {$3$};
  \node[left] at (0, -5) {$-5$};
  \node[left] at (0, -9) {$-9$};
  
  \node[above] at (1, 0) {$1$};
  \node[above] at (-1, 0) {$-1$};
  \node[above] at (1.732, 0) {$\sqrt{3}$};
  \node[above] at (-1.732, 0) {$-\sqrt{3}$};
  \node[below right] at (2.449, 0) {$\sqrt{6}$};
  \node[below left] at (-2.449, 0) {$-\sqrt{6}$};

  % 6. 特徴的なプロット点
  \fill (0, 3) circle (1.2pt);
  \fill (1, -5) circle (1.2pt);
  \fill (-1, -5) circle (1.2pt);
  \fill (1.732, -9) circle (1.2pt);
  \fill (-1.732, -9) circle (1.2pt);

  % 7. 領域 D のラベル
  \node[fill=white, inner sep=1pt, font=\bfseries] at (0, -1.5) {$D$};

\end{tikzpicture}
\end{center}
\caption{領域 $D$ の図示（境界を含まない）}
\end{figure}

\newpage

\end{document}
